\documentclass[10pt,a4paper]{amsart}
\usepackage[margin=19mm]{geometry}
\usepackage{amsmath,amssymb,mathtools}
\usepackage[scr=rsfs,cal=euler]{mathalfa}
\usepackage{xfrac}
\usepackage[unicode,hidelinks,bookmarks=false]{hyperref}
\newcommand{\ie}{i.e.\ }
\newcommand{\cf}{cf.\ }
\newcommand{\resp}{respectively\ }
\newcommand{\Subsection}{\subsection}
\providecommand{\nobreakdash}{\nobreak\hbox{-}}
\setlength{\emergencystretch}{2em}
\multlinegap0pt
\usepackage{amssymb}
\usepackage{enumerate}
\usepackage{xcolor}
\newtheorem{prop}{Proposition}
\newtheorem{thm}{Theorem}
\usepackage{cleveref}
\crefname{prop}{Proposition}{Propositions}
\crefname{thm}{Theorem}{Theorems}
\newcommand{\Cl}{\mc{C}}
\newcommand{\D}{\mathsf{D}}
\newcommand{\G}{\mathsf{G}}
\renewcommand{\H}{\mathsf{H}}
\newcommand\N{\operatorname{N}}
\newcommand\Nrd{\operatorname{\mathsf{Nrd}}}
\newcommand{\Z}{\mathsf{Z}}
\newcommand{\can}{\operatorname{\mathsf{can}}}
\newcommand{\disc}{\mathsf{disc}}
\newcommand\ind{\operatorname{\mathsf{ind}}}
\newcommand{\la}{\langle}
\newcommand{\lla}{\la\!\la}
\newcommand{\mc}{\mathcal}
\newcommand{\mg}[1]{{#1}^{\times}}
\newcommand{\ra}{\rangle}
\newcommand{\rra}{\ra\!\ra}      
\newcommand{\s}{\sigma}
\renewcommand{\setminus}{\smallsetminus}
\newcommand{\sq}[1]{{#1}^{\times 2}}
\title{Non-$R$-trivial proper projective similitudes\\ in type $A_3\equiv D_3$}
\author{M. Archita, Karim Johannes Becher}
\date{Source: arXiv:2605.09110v1 (CC BY 4.0)}
\begin{document}
\maketitle

\noindent\emph{Source and licence.} Non-$R$-trivial proper projective similitudes\\ in type $A_3\equiv D_3$. Version arXiv:2605.09110v1 (CC BY 4.0), distributed by arXiv under the Creative Commons Attribution 4.0 International licence, http://creativecommons.org/licenses/by/4.0/, as recorded in the arXiv licence metadata for this exact version and shown on the abstract page https://arxiv.org/abs/2605.09110v1. Full licence notice: \texttt{licenses/arxiv-2605.09110-CC-BY-4.0.txt}.\par\smallskip

\noindent\textit{Editorial scope.} Counted unit: the article's prop P:3-Pfister-A3-example (source0.tex lines 381--388). The article's complete original proof of it is delivered with it (source0.tex lines 389--394, one \texttt{proof} environment of 6 lines). No mathematics is written, repaired or re-derived: the statement and the proof are the article's own lines, in the article's own notation, and no retained line is substituted or re-flowed.  Retained uncounted as the source-local context the proof needs: lines 312--315; lines 323--336.  Nothing was omitted from any retained span, so the package carries no omission record.  No retained line contains a citation, so the wrapper carries no bibliography and no bibliography entry.  No retained line contains a figure environment, an image-inclusion command or drawing code, so no image is delivered and the package declares no asset dependency.  Editorial formatting only, in addition: an AMS \texttt{amsart} preamble that restates the article's own declarations and macros used by the retained lines, namely the environments \texttt{prop}, \texttt{thm} on separate counters, together with the standard packages those lines need.  The retained environments print in this document as Proposition 1, Theorem 1 in order of appearance, whereas the article prints the same items with the numbering of its own full document.  The complete native source of the article as distributed in the arXiv e-print for this exact version is bundled unchanged as \texttt{sources/200194/source0.tex}.
\begin{prop}\label{D3-ex-construction}
Let $(A,\s)$ be a $K$-algebra with orthogonal involution of degree $6$ and of nontrivial discriminant. Let $C=\Cl(A,\s)$ and $L=\Z(C)$. Assume that $\ind(C)=2$ and $A$ is not split. 
Then there exist $K$-quaternion algebras $Q,Q_1,Q_2$ and an orthogonal involution $\gamma$ on $Q$ such that $L$ is contained in $Q$ and $\gamma$ extends the nontrivial automorphism of $L/K$ and such that $A\sim Q\sim Q_1\otimes Q_2$ and $$(A,\s) \in (Q_1,\can_{Q_1})\otimes (Q_2,\can_{Q_2}) \boxplus (Q,\gamma)\,.$$
\end{prop}

\begin{thm}[Merkurjev] 
\label{T:MerkProp9}
    Let $Q,Q_1,Q_2$ be $K$-quaternion  algebras such that $Q_1\otimes Q_2\sim Q$.
    Let $L/K$ be a quadratic field extension contained in $Q$ and $\gamma$ an orthogonal involution on $Q$ extending the nontrivial automorphism of $L/K$.
    Let $(A,\s)\in (Q_1,\can_{Q_1})\otimes (Q_2,\can_{Q_2})\boxplus (Q,\gamma)$ and $C=\Cl(A,\s)$.
    Then $C$ is a central simple $L$-algebra
%    Assume that $\ind C=2$.
%    Then
with $C\sim (Q_1)_L\simeq (Q_2)_L$, and we have 
$$\G^+(A,\s)  = \N_{L/K}^\ast\cap(\Nrd_{Q_1}^\ast\cdot \Nrd_{Q_2}^\ast) \quad\mbox{ and }\quad
\H(A,\s) = \N^\ast_{L/K}\cap \Nrd^\ast_{Q_k}$$
for $k=1,2$, and in partiuclar
    $${\bf PSim}^+(A,\s)(K)/R\,\,\simeq \,\,\dfrac{\N_{L/K}^\ast\cap(\Nrd_{Q_1}^\ast\cdot \Nrd_{Q_2}^\ast)}{\N^\ast_{L/K}\cap \Nrd^\ast_{Q_k}}\,\,.$$
\end{thm}

\begin{prop}\label{P:3-Pfister-A3-example}
    Let $a,b,c\in\mg{K}$ be such that $-1\in\D_K{\lla a,b\rra}$ and
    $\lla a,b,c\rra$ is anisotropic.
Let $Q_1=(a,b)_K$, $Q_2=(-ac,b)_K$, $Q=(-c,b)_K$ and $L=K(\sqrt{-c})$. 
 Then $$c\in\left(\N_{L/K}^\ast\cap(\Nrd_{Q_1}^\ast\cdot \Nrd_{Q_2}^\ast)\right)\setminus (\Nrd_{Q_1}^\ast\cup\Nrd_{Q_2}^\ast)$$
and there exists an orthogonal involution $\gamma$ on $Q$ with $\disc(\gamma)=-c\sq{K}$.
Furthermore, there exists a $K$-algebra of degree $6$ with orthogonal involution $$(A,\s)\in \left((Q_1,\can_{Q_1})\otimes (Q_2,\can_{Q_2})\right)\boxplus (Q,\gamma)\,,$$ and for any such $(A,\s)$ we have ${\bf PSim}^+(A,\s)(K)/R\neq \{1\}$.
\end{prop}

\begin{proof}
Note that $c=\N_{L/K}(\sqrt{-c})\in\N_{L/K}^\ast$, $-1,-a\in \Nrd_{Q_1}^\ast$ and $ac\in \Nrd_{Q_2}^\ast$. 
Therefore $c\in\N_{L/K}^\ast\cap (\Nrd_{Q_1}^\ast\cdot\Nrd_{Q_2}^\ast)$.
On the other hand, since $\lla a,b,c \rra$ is anisotropic, we have that $c\notin\Nrd_{Q_1}^\ast\cup\Nrd_{Q_2}^\ast$.
The statement now follows by \Cref{D3-ex-construction} and \Cref{T:MerkProp9}.
\end{proof}

\end{document}
